% !TEX root = main.tex
%=====================================
\DeclarePairedDelimiter{\pa}{(}{)}
\DeclarePairedDelimiter{\bra}{[}{]}
\DeclarePairedDelimiter{\set}{\{}{\}}
\DeclarePairedDelimiter{\abs}{\lvert}{\rvert}
\DeclarePairedDelimiter{\norm}{\lVert}{\rVert}
\DeclarePairedDelimiter{\inner}{\langle}{\rangle}
\DeclarePairedDelimiter{\ropen}{[}{)}
\DeclarePairedDelimiter{\lopen}{(}{]}
\NewDocumentCommand{\mfunc}{O{}O{}}{\mathcal{F}_{#1}^{#2}}
\newcommand{\NN}{\mathbb{N}}
\newcommand{\ZZ}{\mathbb{Z}}
\newcommand{\QQ}{\mathbb{Q}}
\newcommand{\RR}{\mathbb{R}}

\newcommand{\PRIMITIVO}[1][]{\mathrm{Prim}_{#1}}
\newcommand{\CERO}{\operatorname{Z}}
\newcommand{\PROYECCION}[2]{\pi_{#2}^{#1}}
\newcommand{\SUCESOR}{\operatorname{S}}
\newcommand{\CPOTENCIA}[1]{2^{#1}}
\newcommand{\CARACTERISTICA}[1]{\chi_{#1}}
\newcommand{\DIVERGE}{\uparrow}
\newcommand{\ACABA}{\downarrow}
\newcommand{\DOMINIO}[1]{\operatorname{dom}\left(#1\right)}
\newcommand{\RANGO}[1]{\operatorname{ran}\left(#1\right)}
\newcommand{\PARCIAL}[1][]{\mathrm{Parcial}_{#1}}
\newcommand{\TOTAL}[1][]{\mathrm{Total}_{#1}}
\newcommand{\INSTRUCCIONES}{\mathrm{Instr}}
\newcommand{\INSTRUCCIONESCOD}{\mathrm{InstrCod}}
\newcommand{\URM}{\mathrm{URM}}
\newcommand{\URMCOMPUTABLE}[1][]{\mathrm{URMComp}_{#1}}
\newcommand{\LONGITUD}{\operatorname{Len}}
\newcommand{\URMCODIGO}{\mathrm{URMCod}}
\newcommand{\MAXREGISTRO}{\mathrm{MaxReg}}
\newcommand{\NUMINSTRUCCIONES}{\mathrm{NumInstr}}
\newcommand{\CONFIGURACIONES}{\mathrm{Config}}
\newcommand{\ARITMETICA}{\mathrm{L}_{2}}
\newcommand{\VARIABLES}[1][]{\mathrm{V}_{#1}}
\newcommand{\MIN}[1]{\min\left(#1\right)}
\newcommand{\SECUENCIAS}{\mathrm{Seq}}
\newcommand{\MAXPOTENCIA}{\mathrm{MaxP}}

\newcommand{\TERMINOS}{\mathrm{Term}}
\newcommand{\ATOMICAS}{\mathrm{Atom}}
\newcommand{\FORMULAS}{\mathrm{Form}}
\newcommand{\ENUNCIADOS}{\mathrm{Sent}}

\newcommand{\INDUCCION}[1]{#1\text{-}\mathrm{IND}}
\newcommand{\COMPRENSION}[1]{#1\text{-}\mathrm{COM}}
\newcommand{\PEANO}{\mathrm{PA}^{-}}
\newcommand{\RCA}{\mathrm{RCA}_{0}}
\newcommand{\ACA}{\mathrm{ACA}_{0}}
\newcommand{\WKL}{\mathrm{WKL}_{0}}
\newcommand{\VARIABLESDE}[1]{\mathrm{Var}\left[#1\right]}
\newcommand{\LIBRESDE}[1]{\mathrm{Free}\left[#1\right]}

\NewDocumentCommand{\PHI}{m o}{% "m" es argumento obligatorio (mandatory), "o" es opcional
	\IfNoValueTF{#2}%
	{\varphi_{#1}}% Si NO se escribe el segundo argumento
	{\varphi_{#1}^{(#2)}}% Si SÍ se escribe el segundo argumento
}